\documentclass[UTF8,fontset=none,11pt,a4paper]{ctexart}
\usepackage{ipara-handbook}
\handbooksetup{NET-01 通信基础与网络性能}{408 · 计算机网络 NET-01}{Signal · Delay · Capacity · BDP}
\providecommand{\tightlist}{\setlength{\itemsep}{0pt}\setlength{\parskip}{0pt}}
\providecommand{\passthrough}[1]{#1}
\begin{document}
\handbooktitle{通信基础与网络性能}{把信息送过有限信道}{408 · NET-01}
\section{0. 本册定位}\label{0-ux672cux518cux5b9aux4f4d}

本 Topic 回答一个问题：一个 bit
要穿过有限带宽、有限传播速度且可能有噪声的真实信道，必须付出什么代价？

本册拥有：source/destination、data/signal/symbol/bit
的边界，传输介质与物理接口，repeater/hub，传输、传播、处理、排队四类时延，throughput、BDP、交换服务与存储转发，Nyquist/Shannon
的 408 模型，以及编码、调制和复用的选择逻辑。

本册不拥有：多站点怎样争用共享介质，见\href{../02_单跳交付_帧_MAC_局域网与交换机/README.md}{单跳交付}；窗口怎样填满管道，见\href{../03_可靠传输_序号_ACK_定时器与滑动窗口/README.md}{可靠传输}；拥塞造成的动态排队反馈，见\href{../07_拥塞_共享资源与反馈控制/README.md}{拥塞控制}。

\section{1.
根本问题：距离不会因为链路变快而消失}\label{1-ux6839ux672cux95eeux9898ux8dddux79bbux4e0dux4f1aux56e0ux4e3aux94feux8defux53d8ux5febux800cux6d88ux5931}

通信不是``把文件瞬移到另一台机器''，而是连续完成三件事：

\begin{enumerate}
\def\labelenumi{\arabic{enumi}.}
\tightlist
\item
  把信息表示成接收方能够区分的物理状态；
\item
  用有限速率把这些状态注入介质；
\item
  等待扰动沿有限距离传播到另一端。
\end{enumerate}

因此网络性能的第一母模型是：

\[
\boxed{
\text{Information}
\to \text{Symbol}
\to \text{Signal}
\to \text{Link}
\to \text{Arrival}
}
\]

箭头表示信息逐步获得物理表示并在空间中传播，不是协议分层。它导出两个不能混合的时钟：

\[
\boxed{T_{tx}=\frac{L}{R}}
\qquad
\boxed{T_{prop}=\frac{d}{v}}
\]

\begin{itemize}
\tightlist
\item
  \(T_{tx}\) 决定``最后一个 bit 何时离开发送端''，由数据长度 \(L\)
  和发送速率 \(R\) 决定；
\item
  \(T_{prop}\) 决定``某个 bit 何时到达接收端''，由距离 \(d\) 和传播速度
  \(v\) 决定。
\end{itemize}

提高带宽会缩短发送时延，却不会消除传播时延。这是后续窗口、流水线和拥塞机制出现的物理起点。

\section{2.
对象不等于表示}\label{2-ux5bf9ux8c61ux4e0dux7b49ux4e8eux8868ux793a}

一次通信的最小物理链不是“发送端直接把 bit 给接收端”，而是：

\[
\boxed{
\text{Source}\to\text{Transmitter}\to\text{Channel}\to
\text{Receiver}\to\text{Destination}}
\]

Source 产生信息，transmitter 把信息编码/调制成适合介质的信号，channel 传播并叠加衰减、失真与噪声，receiver 从观测信号恢复数据，destination 消费数据。信源/信宿是信息语义的两端；发送器/接收器是物理表示转换的两端，不能因为它们常在同一设备中就合并概念。

\begin{longtable}[]{@{}lclll@{}}
\toprule\noalign{}
概念 A & ≠ & 概念 B & 真正区别与题目信号 & 混淆后果 \\
\midrule\noalign{}
\endhead
\bottomrule\noalign{}
\endlastfoot
Data & ≠ & Signal & Data 是要表达的信息；signal 是介质中的物理载体 &
把``数字数据''误认为只能用数字信号传输 \\
Bit & ≠ & Symbol & bit 是信息单位；symbol 是一次发送的可区分物理状态 &
漏掉每码元可承载 \(\log_2 M\) bit \\
Bit rate & ≠ & Baud rate & 前者数 bit/s，后者数 symbol/s & 曼彻斯特、QAM
等题中倍数关系算反 \\
Bandwidth & ≠ & Throughput & 前者是链路能力上限；后者是实际交付速率 &
忽略瓶颈、协议等待和拥塞 \\
Transmission & ≠ & Propagation & 前者是注入链路；后者是信号移动 &
用距离计算发送时延或用包长计算传播时延 \\
\end{longtable}

\section{3. 传输介质与物理接口：信号能走，不等于双方会通信}

介质决定信号传播的物理约束，却不单独决定 bit rate、协议或可靠性。比较介质时使用四个维度：可用频带与衰减、抗干扰能力、部署距离与成本、是否易被共享或窃听。

\begin{handbooktable}{L{2.2cm} Y Y Y}
\toprule
介质 & 物理直觉 & 主要优势 & 主要边界\\
\midrule
双绞线 & 两根绝缘铜线绞合，使外界干扰在两线中尽量共模抵消 & 成本低、布线方便 & 衰减和串扰限制距离/速率，类别与题设相关\\
同轴电缆 & 中心导体与同轴屏蔽层形成受控传输结构 & 屏蔽和带宽通常优于普通双绞线 & 较粗重，现代 LAN 主干中不再是唯一选择\\
光纤 & 光在纤芯/包层结构中传播 & 带宽高、衰减低、抗电磁干扰 & 光电转换与施工成本；单模/多模边界不同\\
无线 & 电磁波在开放空间传播 & 移动性、广播可达、部署灵活 & 干扰、衰落、共享竞争和安全暴露\\
\bottomrule
\end{handbooktable}

“光纤一定传播得更快”不是稳定结论。介质的传播速度与链路发送速率是两个维度；光纤常见优势是可用带宽、衰减和抗干扰，而不是把传播时延变为零。

物理接口还必须让两端就四类约束达成一致：

\begin{itemize}
\tightlist
\item \textbf{机械特性：}连接器形状、引脚数量和物理尺寸；
\item \textbf{电气特性：}电平、阻抗、速率和信号持续时间；
\item \textbf{功能特性：}每条线或信号的用途；
\item \textbf{规程特性：}信号按什么先后次序出现、何时有效。
\end{itemize}

只有介质而没有接口契约，双方即使物理相连也不能正确解释信号。

\subsection{Repeater 与 Hub：只延伸 bit 的可达范围}

Repeater 在物理层接收衰减/失真的信号并重新整形、定时、发送，不读取 MAC 地址，也不隔离上层广播。Hub 是多端口 repeater：从一个端口观察到的 bit 被再生到其他端口，所有站共享带宽与冲突域。它没有 switch 的 source learning、MAC table 和定向转发状态。

若每个码元有 \(M\) 个可可靠区分的状态，则理想映射下：

\[
R_b=R_s\log_2 M.
\]

增加 \(M\) 能让一个码元携带更多
bit，但相邻状态更难区分，所需信噪比也更高。它不是免费的速率提升。

\section{4.
容量上限怎样生成}\label{3-ux5bb9ux91cfux4e0aux9650ux600eux6837ux751fux6210}

\subsection{4.1
无噪声仍受带宽限制}\label{31-ux65e0ux566aux58f0ux4ecdux53d7ux5e26ux5bbdux9650ux5236}

朴素想法是无限提高码元变化速度。失败原因是有限带宽会使相邻码元的波形互相干扰。408
使用 Nyquist 模型表达无噪声低通信道的上限：

\[
R_{s,\max}=2W,
\qquad
R_{b,\max}=2W\log_2M.
\]

这里 \(W\) 是 Hz，\(M\)
是码元状态数。第一式约束每秒可区分多少码元，第二式再乘每码元的信息量。

\subsection{4.2
有噪声时不能无限增加状态数}\label{32-ux6709ux566aux58f0ux65f6ux4e0dux80fdux65e0ux9650ux589eux52a0ux72b6ux6001ux6570}

若只看 Nyquist，增大 \(M\) 似乎能无限提高 bit
rate。现实噪声会让密集状态不可区分。Shannon
容量给出有噪声信道的可靠通信极限：

\[
C=W\log_2(1+\mathrm{SNR}).
\]

公式中的 SNR 是功率比；若题目给分贝：

\[
\mathrm{SNR}_{dB}=10\log_{10}\mathrm{SNR}.
\]

Nyquist 回答``带宽允许多快地摆放可区分码元''，Shannon
回答``噪声存在时最多能可靠携带多少信息''。题目同时给出两组限制时，系统必须同时满足，不能任选一个公式。

\section{5.
编码、调制与复用解决不同问题}\label{4-ux7f16ux7801ux8c03ux5236ux4e0eux590dux7528ux89e3ux51b3ux4e0dux540cux95eeux9898}

\subsection{5.1 编码：怎样把 bit
变成可同步的基带波形}\label{41-ux7f16ux7801ux600eux6837ux628a-bit-ux53d8ux6210ux53efux540cux6b65ux7684ux57faux5e26ux6ce2ux5f62}

线路编码必须同时考虑：接收方能否恢复时钟、是否存在难以传输的直流分量、付出多少额外码元。

\begin{itemize}
\tightlist
\item
  NRZ 简单，但长串同值缺少跳变；
\item
  Manchester 用位中跳变携带时钟，易同步但码元开销大；
\item
  4B/5B 先限制长串零，再配合 NRZI，在同步能力和效率之间折中。
\end{itemize}

编码不是可靠传输。它改善物理可判决性，却不替代 CRC、ACK 或重传。

\subsection{5.2
调制：怎样把信息装到适合信道的载波上}\label{42-ux8c03ux5236ux600eux6837ux628aux4fe1ux606fux88c5ux5230ux9002ux5408ux4fe1ux9053ux7684ux8f7dux6ce2ux4e0a}

载波可写作 \(A\cos(2\pi ft+\phi)\)。ASK、FSK、PSK
分别改变振幅、频率或相位；QAM
联合改变振幅和相位。调制阶数越高，每码元携带的 bit
越多，但判决间隔越小、对 SNR 越敏感。

\subsection{5.3
复用：怎样让多个信息流共享一条介质}\label{43-ux590dux7528ux600eux6837ux8ba9ux591aux4e2aux4fe1ux606fux6d41ux5171ux4eabux4e00ux6761ux4ecbux8d28}

FDM、TDM、WDM、CDM
分别在频率、时间、光波长或码空间上分离使用者。复用先划分资源；MAC
则在共享资源未被预先完全划分时决定``现在轮到谁''。两者在概念上相邻，但责任不同。

\begin{handbooktable}{L{1.8cm} Y Y Y}
\toprule
方式 & 正交/分离轴 & 必须支付的控制成本 & 典型计算接口\\
\midrule
FDM & 不重叠频带 & guard band、滤波 & 总带宽含各子带与保护带\\
TDM & 周期性 time slot & frame/slot synchronization & 每轮槽数、每槽 bit 数、frame rate 决定线路 bit rate\\
WDM & 光纤中的不同 wavelength & 光器件与通道间隔 & 可近似视为光域 FDM\\
CDM & 近似正交 code sequence & chip rate、同步、相关检测 & 用内积/相关值分离用户\\
\bottomrule
\end{handbooktable}

同步 TDM 即使某输入暂时无数据，也可能保留其固定 slot；统计 TDM 可动态分配槽位提高突发业务利用率，却需地址/调度与缓冲。若 $n$ 路输入每路每秒产生 $f$ 个样本、每样本 $b$ bit，且一帧各取一个样本，忽略额外开销时线路速率至少为 $nfb$ bit/s；题设给同步位或 framing overhead 时必须另加。CDM 的“同时同频”不等于无干扰，它依赖码序列相关性与接收同步。

\section{6. 交换服务：状态何时建立，节点按多大粒度等待}

交换方式可以用两个问题生成：传输前是否建立沿途状态？中间节点必须收齐多大对象才能继续？

\begin{handbooktable}{L{2.2cm} Y Y Y}
\toprule
模型 & 状态与资源 & 中间节点动作 & 主要代价\\
\midrule
电路交换 & 传输前建立端到端连接并预留资源 & 建立后连续转送 bit 流 & setup 时间；空闲时预留能力也不能被他人使用\\
报文交换 & 不预留端到端电路 & 收齐整个 message 后存储、排队、选择下一跳 & 大报文要求大缓冲，并阻塞后续短报文\\
分组交换 & 通常不预留完整物理电路；把 message 切成 packet & 按 packet 存储转发，可形成流水线 & 每包首部、排队、乱序/丢失与重组责任\\
\bottomrule
\end{handbooktable}

报文交换与分组交换都可 store-and-forward，区别是缓存/调度粒度。packet 化不能消除总数据发送量，但允许第一包在整份 message 尚未发送完时进入下一跳，也让多个流细粒度交错。

\subsection{数据报与虚电路：同为分组交换，状态位置不同}

\begin{handbooktable}{L{2.2cm} Y Y Y}
\toprule
模型 & 建立阶段 & packet 携带/网络保存 & 故障与顺序直觉\\
\midrule
Datagram & 无连接建立要求 & 每包携带完整目的标识，router 按当前转发表独立处理 & 不同包可走不同路径、可能乱序；路由变化可影响后续包\\
Virtual circuit & 先建立逻辑路径/表项 & 包携带较短 VC identifier；沿途设备保存映射状态 & 同一 VC 通常沿既定逻辑路径；沿途状态故障需要恢复/重建\\
\bottomrule
\end{handbooktable}

虚电路是网络中的逻辑状态，不等于物理专线，也不必等于预留固定带宽；数据报无连接也不等于“网络没有任何状态”，router 仍保存 forwarding state。IP 的具体数据报转发由 NET04 Own，本节只拥有交换服务模型。

\section{7.
一个分组经过链路的一生}\label{5-ux4e00ux4e2aux5206ux7ec4ux7ecfux8fc7ux94feux8defux7684ux4e00ux751f}

设长度为 \(L\) 的 packet 依次经过 \(N\)
条等速率链路，每个中间交换节点采用 store-and-forward，暂忽略排队与处理：

\begin{lstlisting}
源端序列化第 1 个 packet
-> 第 1 条链路传播
-> 中间节点收齐后才能向第 2 条链路序列化
-> ...
-> 目的端收齐
\end{lstlisting}

单个 packet 的最早完整到达时间为：

\[
T=\sum_{i=1}^{N}\left(\frac{L}{R_i}+\frac{d_i}{v_i}\right).
\]

若有 \(P\) 个等长 packet、\(N\)
条等速率链路且可形成稳定流水线，则忽略其他时延时：

\[
T=(N+P-1)\frac{L}{R}+\sum_{i=1}^{N}T_{prop,i}.
\]

\(N+P-1\) 来自``第一个 packet 填满 \(N\) 级流水线，余下 \(P-1\)
个各增加一个发送周期''，不能机械背成适用于任意速率和任意交换方式的公式。

\section{8.
四类时延与瓶颈}\label{6-ux56dbux7c7bux65f6ux5ef6ux4e0eux74f6ux9888}

每个节点的总时延拆为：

\[
d_{nodal}=d_{proc}+d_{queue}+d_{trans}+d_{prop}.
\]

\begin{itemize}
\tightlist
\item
  processing：解析首部、检错和查表；
\item
  queueing：等待输出资源，随 offered load 动态变化；
\item
  transmission：把所有 bit 推上当前链路；
\item
  propagation：信号穿过当前介质。
\end{itemize}

端到端 throughput 由路径上持续服务能力最弱的环节限制：

\[
R_{throughput}\le \min_i R_i.
\]

但``最小链路速率''等于实际吞吐量只在发送方、接收方、协议窗口和拥塞均不形成更小瓶颈时成立。

\section{9.
BDP：网络中同时在途的状态预算}\label{7-bdpux7f51ux7edcux4e2dux540cux65f6ux5728ux9014ux7684ux72b6ux6001ux9884ux7b97}

带宽时延积不是``又一种时延''，而是一定传播时间内链路可注入的数据量：

\[
BDP=R\times T_{prop}.
\]

若讨论往返反馈闭环，常使用 \(R\times RTT\)
估算填满路径所需的未确认数据量。它将物理 Topic 与可靠传输、TCP
流控和拥塞控制连接起来：窗口小于在途容量时，发送方会等待反馈；窗口过大又可能积压队列。

\section{10.
正确性、代价与边界}\label{8-ux6b63ux786eux6027ux4ee3ux4ef7ux4e0eux8fb9ux754c}

\begin{longtable}[]{@{}llll@{}}
\toprule\noalign{}
机制 & 换来的能力 & 主要代价 & 失效或边界 \\
\midrule\noalign{}
\endhead
\bottomrule\noalign{}
\endlastfoot
更高阶调制 & 每码元更多 bit & 更高 SNR 与实现复杂度 &
噪声使星座点不可区分 \\
分组交换 & 突发业务统计复用 & 排队、抖动、丢包 &
高负载时延可能急剧增加 \\
Store-and-forward & 完整检错和逐跳转发 & 每跳增加一个 packet 发送时延 &
cut-through 模型不能直接套用 \\
复用 & 多流共享昂贵介质 & 划分、同步或码设计开销 &
资源空闲时可能利用不足 \\
更高速率链路 & 缩短 transmission delay & 成本与物理实现要求 &
不改变距离造成的 propagation delay \\
\end{longtable}
\section{11. 408 流程覆盖：选择、交换与时间线}

\subsection{编码与调制选择不是名词配对}

题目给定数据、信道和目标速率时，按以下顺序推进：
\[
\boxed{\text{Data type}\to\text{Baseband/Passband}\to\text{Symbol states}\to\text{Rate limits}\to\text{Feasibility}}
\]
先把 bit rate 换成 symbol rate，再同时检查 Nyquist 与 Shannon 上限。若候选方案超过任一上限，结论是该表示在给定条件下不可行，而不是继续提高码元状态数。

\subsection{电路、报文与分组交换的事件账本}

电路交换先经历资源建立，再持续占用预留能力，最后释放；报文交换逐份完整 message 排队；分组交换逐包排队、逐跳存储转发：
\[
\text{Setup}\to\text{Reserved transfer}\to\text{Release}
\]
\[
\text{Whole message arrival}\to\text{Store}\to\text{Forward}
\qquad\text{vs}\qquad
\text{Packet arrival}\to\text{Queue}\to\text{Serialize}\to\text{Propagate}.
\]
因此计算端到端时间时，必须先问是否存在 setup cost、是否 store-and-forward、多个 packet 是否能形成流水线，以及各链路速率是否相同。

\begin{examplebox}{三跳两包时间线}
两包经过三条等速率 store-and-forward 链路。第一包要经历三个发送周期才能到达；此后第二包已在上一跳流水推进，只再增加一个发送周期，所以发送项为 $(3+2-1)L/R$，不是 $2\times3L/R$。传播、处理和排队仍按题设另加。
\end{examplebox}

\section{12. 做题调用协议}\label{9-ux505aux9898ux8c03ux7528ux534fux8bae}

\begin{enumerate}
\def\labelenumi{\arabic{enumi}.}
\tightlist
\item
  标出对象：bit、symbol、frame、packet 还是整份数据；
\item
  统一单位：bit/Byte，Hz，bit/s，s；
\item
  画时间线，分别标注 serialization、propagation、processing、queueing；
\item
  判断是单链路、逐跳 store-and-forward 还是多 packet 流水线；
\item
  交换题先写是否 setup、节点按 message 还是 packet 缓存、网络是否保存 per-connection state；
\item
  容量题先判断 Nyquist、Shannon 或二者共同约束；
\item
  复用题先识别资源轴，再把 guard/sync/address/code 等开销计入总预算；
\item
  吞吐量题列出所有候选瓶颈，再取最小；
\item
  用数量级和量纲独立检查结果。
\end{enumerate}

\section{13.
贯穿母例：高速远距离链路为什么仍需要窗口}\label{10-ux8d2fux7a7fux6bcdux4f8bux9ad8ux901fux8fdcux8dddux79bbux94feux8defux4e3aux4ec0ux4e48ux4ecdux9700ux8981ux7a97ux53e3}

一条 \(100\,\mathrm{Mb/s}\) 链路，单向传播时延 \(20\,\mathrm{ms}\)，发送
\(1000\,\mathrm{B}\) frame。

\[
T_{tx}=\frac{8000}{10^8}=80\,\mu s,
\qquad
RTT\approx40\,ms.
\]

停止等待时，发送方工作 \(80\,\mu s\) 后大约等待
\(40\,ms\)。问题不是链路``慢''，而是反馈尚未返回。对应的往返在途容量约为：

\[
R\times RTT=4\times10^6\,bit=500\,kB.
\]

所以需要允许大量未确认数据同时在途。这个例子在\href{../03_可靠传输_序号_ACK_定时器与滑动窗口/README.md}{可靠传输}中会生成
sliding
window，在\href{../06_传输层_端点_UDP与TCP状态机/README.md}{TCP}中分别受到
\passthrough{\lstinline!rwnd!} 和 \passthrough{\lstinline!cwnd!} 约束。

\section{14. 高频 First
Divergence}\label{11-ux9ad8ux9891-first-divergence}

\begin{itemize}
\tightlist
\item
  看到``100 Mb/s''就用距离计算：把 transmission 与 propagation 混合；
\item
  看到带宽和 SNR 只算一个上限：没有求共同可行域；
\item
  把 packet 数乘在整条路径总时延上：没有识别流水线重叠；
\item
  把 BDP 当成已发送文件总量：没有说明它描述哪段路径、哪个时间尺度；
\item
  只找最慢链路：漏掉窗口、发送源或拥塞可能成为更小瓶颈。
\item
  把虚电路等同于物理电路：没有区分逻辑转发表状态与资源预留；
\item
  说 hub 根据 MAC 定向发送：把物理层 bit 再生错写成链路层转发。
\end{itemize}

\section{15.
一页压缩与复原问题}\label{12-ux4e00ux9875ux538bux7f29ux4e0eux590dux539fux95eeux9898}

\[
\boxed{
\text{Represent}
\to \text{Inject at finite rate}
\to \text{Propagate over distance}
\to \text{Queue at shared resources}
\to \text{Arrive under a bottleneck}
}
\]

复原本册时回答：

\begin{enumerate}
\def\labelenumi{\arabic{enumi}.}
\tightlist
\item
  为什么提高带宽不一定缩短 RTT？
\item
  Nyquist 与 Shannon 分别限制什么？
\item
  为什么 packet 化会增加单包开销，却改善突发业务共享？
\item
  多跳多 packet 总时延中的流水线项怎样生成？
\item
  BDP 为什么会自然导出窗口？
\item
  数据报与虚电路把状态分别放在哪里？
\item
  Repeater、hub 与 switch 读取的信息为何不同？
\end{enumerate}

\section{16.
来源与校正说明}\label{13-ux6765ux6e90ux4e0eux6821ux6b63ux8bf4ux660e}

\begin{itemize}
\tightlist
\item
  归档笔记《网络概述》《物理层-编码与调制》《公式汇总》提供考点范围和旧例题；
\item
  旧笔记中的公式已按对象、作用域和前提重新组织，未保留``看到关键词套公式''的结构；
\item
  Nyquist/Shannon 在本册按 408
  教材模型使用，工程通信中的脉冲成形、编码增益和具体信道模型不在本册展开。
\item
  Ethernet 介质与 repeater 的工程边界参考 IEEE 802.3-2022；本册只保留 408 所需的介质、PHY 与物理设备模型。
\end{itemize}
\end{document}
